\documentclass{article} % This command is used to set the type of
% document you are working on such as an article, book, or presenation

\usepackage{
  amsmath,      % Math Environments
  amssymb,      % Extended Symbols
  enumerate,        % Enumerate Environments
  graphicx,      % Include Images
  lastpage,      % Reference Lastpage
  multicol,      % Use Multi-columns
  multirow      % Use Multi-rows
}

\usepackage{geometry} % This package allows the editing of the page layout
\usepackage{amsmath}  % This package allows the use of a large range
% of mathematical formula, commands, and symbols
\usepackage{graphicx}  % This package allows the importing of images

\newcommand{\question}[2][]{
  \begin{flushleft}
    \textbf{Question #1}: \textit{#2}

\end{flushleft}}
\newcommand{\sol}{\textbf{Solution}:} %Use if you want a boldface solution line
\newcommand{\maketitletwo}[2][]{
  \begin{center}
    \Large{\textbf{Assignment #1}

    N. Ohayon Rozanes} % Name of course here
    \vspace{5pt}

    \normalsize{N. Ohayon Rozanes  % Your name here

    \today}        % Change to due date if preferred
    \vspace{15pt}

\end{center}}
\begin{document}
\maketitletwo[5]  % Optional argument is assignment number
%Keep a blank space between maketitletwo and \question[1]

\question[1]{}
\begin{enumerate}
  \item
    \[
      \int_1^2 \delta(x)(x-2)dx
    \]
    Clearly, since the bounds of integration do not include the point
    at which $\delta(x)$ is concentrated, $0$, the integrand is $0$
    at all points on the interval $[1,2]$ so the integral evaluates $0$.
  \item
    \[
      \int_0^3 \delta_1(x)e^x dx
    \]
    Since $1 \in [0,3]$ then the $\delta_1$ function serves to
    "sample" the other function in the integrand $e^x$ at $x = 1$, so
    the integral evaluates to $e^1$
\end{enumerate}

\question[2]{}

\begin{enumerate}
  \item $h(x) = x\delta(x-1) = x\delta_1(x)$
    However, since $\delta_1$ is only non zero at $x=1$, we can
    further simplify $h$ to be
    \[
      h(x) = \delta_1(x)
    \]
    Which we interpret as defining the functional
    \[
      L_h[u] = \langle h, u \rangle = \langle \delta_1(x), u \rangle
      = u(1)
    \]
  \item
    $h(x) = 3\delta_1(x) - 3x\delta_{-1}(x) = 3(\delta_1(x) - x\delta_{-1}(x))$
    As with the previous question, $\delta_{-1}(x)$ is only non-zero
    at $x=-1$ so distributionally $x\delta_{-1}(x) = -\delta_{-1}(x)$

    So we simplify the $h$ as
    \[
      h(x) = 3(\delta_1(x) + \delta_{-1}(x))
    \]
    On a test function $u$
    \[
      L_h[u] =
      \begin{cases}
        3u(1) & x = 1\\
        3u(-1) & x = -1
      \end{cases}
    \]
\end{enumerate}

\question[3]{}
\begin{enumerate}
  \item We first rewrite the function in terms of $\sigma$ distributions
    \begin{align*}
      f(x) = (\sigma(x+1) - \sigma(x))(x+1) + (\sigma(x) - \sigma(x-1))(1-x)
    \end{align*}
    then
    \[
      f'(x) = (\delta(x+1) - \delta(x))(x+1) + (\sigma(x+1) -
      \sigma(x)) + (\delta(x) - \delta(x-1))(1-x) - (\sigma(x) - \sigma(x-1))
    \]
    The $\delta$ terms vanish, so we can rewrite it more simply as:
    \[
      f'(x) =
      \sigma(x+1) - 2\sigma(x) + \sigma(x-1)
      =
      \begin{cases}
        0 & x < -1 \\
        1 & -1 < x < 0 \\
        -1 & 0 < x < 1 \\
        0 & x > 1
      \end{cases}
    \]
  \item Of course, $k$ can be rewritten as:
    \[
      k(x) =
      \begin{cases}
        0 & x<-2\\
        -x &-2<x<0 \\
        x & 0<x<2\\
        0 & 2<x
      \end{cases}
    \]
    Clearly, there are jump discontinuities at $x=-2$,$x=2$, so
    the piecewise derivative is
    \[
      \begin{cases}
        0 & x<-2\\
        -1 &-2<x<0 \\
        1 & 0<x<2\\
        0 & 2<x
      \end{cases}
    \]
    with $\delta$ functions at the jump discontinuities of maginitude
    $2$, then:
    \[
      k'(x) = 2\delta_{-1}(x) + 2\delta_1(x) +
      \begin{cases}
        0 & x<-2\\
        -1 &-2<x<0 \\
        1 & 0<x<2\\
        0 & 2<x
      \end{cases}
    \]
\end{enumerate}
\end{document}
